\section{Evidence-Traceable Design Rules}
The following rules summarize the cross-cutting synthesis: (1) do not infer ship-scale drag benefit from settlement inhibition alone; (2) report the fouling stage and organism; (3) treat roughness and adhesion after ageing as primary endpoints; (4) separate activity from release and toxicity; (5) compare cleaning force with coating damage and waste generation; (6) use NR rather than filling gaps with assumed values; and (7) assess hulls, nets, membranes, and other assets within their own operating boundary.

For operational cleaning, define the safe cleaning window as
\begin{equation}
\tau_{\mathrm{adhesion}} < \tau_{\mathrm{clean}} < \tau_{\mathrm{damage}},
\end{equation}
where the applied cleaning stress must exceed the target organism's adhesion threshold while remaining below the coating's damage or yield threshold. The local waterjet study reports effective cleaning over approximately $80$--$530$ Pa wall shear under its tested conditions, with higher low-standoff test conditions reaching approximately $1.3$ kPa; these values are not universal limits. High-shear transit hulls and static marine installations require different operating envelopes because hydrodynamic exposure, fouling pressure, access, and cleaning frequency differ. Coating-specific adhesion and wear measurements are required before selecting a trigger or force.

\begin{table*}[t]
\centering
\caption{Taxonomy of sustainable marine biofouling-control strategies. NR denotes not reported in the available corpus.}
\label{tab:taxonomy}
\small
\begin{tabularx}{\textwidth}{l l l X X}
\toprule
Control mode & Representative systems & Target stage & Main mechanism & Principal limitation \\
\midrule
Antifouling & Metal oxides, biocides, bio-derived compounds & Settlement and early biofilm & Toxicity, reactive species, deterrence, or chemistry-mediated inhibition & Release and non-target effects; efficacy is condition-dependent \\
Fouling-release & Silicone, silicone composites, low-modulus surfaces & Established soft and hard fouling & Low adhesion and interfacial fracture under motion or cleaning & Requires retained smoothness, elasticity, and suitable operating conditions \\
Anti-adhesion & Hydrated, zwitterionic, heparin, hydrogel, amphiphilic surfaces & Bacteria, diatoms, larvae & Hydration layer, steric barrier, or reduced interfacial interaction & Ageing, contamination, and mechanical durability \\
Self-cleaning & Photocatalytic, UV, ciliary, lubricant-infused surfaces & Deposits and early fouling & Stimulus-driven degradation, movement, or low-shear release & Energy, optical access, lubricant loss, or scale-up constraints \\
Active maintenance & Grooming, waterjet, ultrasonic, robot, laser & Established fouling and niche areas & Mechanical or energy-assisted removal & Coating wear, waste capture, access, and operational cost \\
\bottomrule
\end{tabularx}
\end{table*}

\begin{table*}[t]
\centering
\caption{Environmental and life-cycle risk pathways identified in the corpus. Values not consistently reported are marked NR.}
\label{tab:risk}
\small
\begin{tabularx}{\textwidth}{l l X X}
\toprule
Source or activity & Release pathway & Reported concern & Mitigation opportunity \\
\midrule
Copper, zinc, silver, and oxide coatings & Dissolution, particle detachment, or cleaning loss & Exposure of non-target organisms and uncertain chronic effects & Retention testing, release measurement, and scenario-specific ecotoxicology \\
Nanocomposite coatings & Nanoparticle or matrix-fragment release during ageing and abrasion & Environmental fate and bioavailability are incompletely resolved & Measure release alongside antifouling performance and wear \\
Hydroblasting and in-water cleaning & Paint particles, microplastics, metals, and fouling biomass in effluent & Toxicity and dispersal to receiving waters & Capture, filtration, treatment, and regulated discharge \\
Biocide-free polymer surfaces & Polymer production, wear, replacement, and disposal & Biocide-free does not imply zero life-cycle impact & Functional-unit LCA including service life and maintenance \\
Active UV, laser, and robotic systems & Energy use and operational waste & Energy, safety, access, and containment constraints & Condition-based deployment and asset-specific energy accounting \\
\bottomrule
\end{tabularx}
\end{table*}

\begin{table*}[t]
\centering
\caption{Proposed tiered release-testing protocol for nanocomposite and photocatalytic marine coatings. The protocol is a methodological recommendation synthesized from the reviewed release, cleaning, and coating studies; it is not a reported dataset.}
\label{tab:release_protocol}
\small
\begin{tabularx}{\textwidth}{l l X X X}
\toprule
Tier & Exposure design & Fractions and measurements & Required controls & Decision output \\
\midrule
1: baseline immersion & Closed-bottle natural or artificial seawater; report exposed area, coating mass, water-to-area ratio, salinity, pH, temperature, dissolved oxygen, light, and sampling times & Dissolved, colloidal, and particulate fractions; ICP-MS/OES for elements; particle number and size by single-particle ICP-MS, microscopy, or validated orthogonal method; coating mass balance & Field and procedural blanks, unfilled matrix, neat substrate, duplicate samples, matrix spikes, filtration recovery, and particle-loss check & Area-normalized concentration--time profile and fraction-specific mass balance \\
2: service stress & Add representative flow, calibrated abrasion or waterjet exposure, wet--dry cycling, UV, and repeated cleaning; document stress energy and exposure history & Repeat Tier 1 measurements plus mass loss, roughness, wettability, adhesion, coating integrity, and antifouling endpoint before and after stress & Untreated coating, stress-only blank, positive release control, and pre/post surface characterization & Retained performance versus release and wear under a defined service scenario \\
3: mesocosm or field & Matched panels or coupons in a controlled mesocosm or field deployment; recover water, sediment or biofilm, and detached solids & Dissolved and particulate release, organism/community response, fouling, surface condition, and environmental chemistry; record hydrodynamic and cleaning history & Spatial and temporal replicates, reference panels, recovery efficiency, chain of custody, and receiving-water background & Asset-relevant exposure estimate with uncertainty and ecological context \\
Ecotoxicity module & Use receiving-environment-relevant primary producer or algae, invertebrate or larval, and microbial endpoints; separate dissolved-ion and particle-only exposures where feasible & Concentration--response and time-response data for matched fractions; report measured rather than nominal exposure where possible & Solvent, seawater, dissolved-ion, particle-only, and positive toxicity controls; blinded sample coding where practical & Hazard characterization linked to the same fractions and concentrations measured in release testing \\
\bottomrule
\end{tabularx}
\end{table*}

\begin{table*}[t]
\centering
\caption{Mapping of the proposed release-testing workflow to recognized standards. The workflow is complementary where it adds marine coating-specific fractionation or service-stress information; it does not supersede the cited methods. Current editions and matrix-specific applicability must be verified before laboratory use.}
\label{tab:standards_mapping}
\small
\begin{tabularx}{\textwidth}{l l X X}
\toprule
Protocol element & Recognized method & Use in this review & Complementary versus superseding elements \\
\midrule
Dissolved biocide/metal release & ISO 15181-1 and the ISO 15181 Parts 1--6 family & Extract and quantify released biocides from antifouling paint under the applicable part-specific procedure & Complementary fractionation of dissolved, colloidal, and particulate material; does not supersede ISO 15181 analyte-specific procedures \\
Seawater matrix & ASTM D1141 & Prepare substitute ocean water for controlled baseline immersion and release comparisons & Complementary natural-seawater and salinity controls; does not replace ASTM D1141 composition or reporting requirements \\
Hydrodynamic/service wear & ASTM D4060; ISO 7784 abrasion-resistance family & Apply controlled abrasion and quantify coating mass loss, roughness, and post-stress release/performance & Complementary waterjet and flow exposure represents service shear; neither abrasion method supersedes coating-specific hydrodynamic validation \\
Tiered algal ecotoxicity & OECD TG 201 & Test growth inhibition of a primary producer/algal endpoint using matched release fractions where matrix adaptation is justified & Complementary measured-exposure and particle/dissolved-ion controls; not a substitute for marine-specific validation without documenting adaptation \\
Tiered invertebrate ecotoxicity & OECD TG 202 & Test acute immobilization of a Daphnia endpoint as a standardized comparative toxicity screen & Complementary marine invertebrate testing is still needed; the freshwater method does not supersede receiving-environment tests \\
Tiered fish ecotoxicity & OECD TG 203 & Test acute fish toxicity for a matched release fraction when scientifically and ethically justified & Complementary marine species and chronic tests; does not establish lifecycle safety \\
Wastewater-treatment relevance & OECD TG 209 & Screen activated-sludge respiration inhibition for cleaning effluent or treatment residuals & Complementary wastewater-treatment endpoint; it does not replace receiving-water ecotoxicity or paint-release measurements \\
\bottomrule
\end{tabularx}
\end{table*}

\begin{table*}[t]
\centering
\caption{Condition-specific quantitative anchors from the curated corpus. These values are not pooled ranges; NR denotes that the requested endpoint was not reported or was not comparable.}
\label{tab:quantitative_anchors}
\small
\begin{tabularx}{\textwidth}{l l l X X}
\toprule
Endpoint & Reported value & System and condition & Interpretation & Source \\
\midrule
Biofilm thickness & $33.6\pm14.8~\mu$m for groomed panels; ungroomed significantly thicker & Fouling-release panels with grooming comparison & Direct panel measurement; not a universal LoF threshold & Hunsucker et al. \\
Frictional drag force & $0.75\pm0.09$ N groomed vs $1.09\pm0.06$ N ungroomed & Matched groomed/ungroomed biofilms in a hydrodynamic test chamber & Relative comparison under the reported panel and flow conditions & Hunsucker et al. \\
Biofilm shaft-power context & Approximately $1$--$18\%$ in literature summarized by the source study & Ship-scale context, not a single experiment in the present corpus & Contextual expectation only; requires vessel-specific ISO 19030/ITTC validation & Snowdon et al. \\
Cleaning wall shear & Approximately $80$--$530$ Pa for visually effective cleaning in tested conditions & Immersed waterjet, coating panels & Condition-specific cleaning force range; coating compatibility must be checked & Oliveira and Granhag \\
Cleaning upper test condition & Approximately $1.3$ kPa wall shear and $0.17$ MPa stagnation pressure at low standoff & Later monthly-cleaning events in the same study & Upper test condition, not a universal operating limit & Oliveira and Granhag \\
Copper release concentration & $178.7\pm43.7$ mg/L estuarine; $33.1\pm7.3$ mg/L seawater after 120 days & Batch leachate from a non-nano copper coating & Concentration is not an area-normalized field release rate & Miller et al. \\
Copper release rate & Approximately $13~\mu$g~cm$^{-2}$~day$^{-1}$ over 364 days & Coated panels under the reported immersion/cleaning conditions & Equivalent release estimate; formulation and cleaning dependent & Oliveira and Granhag \\
Lubricant viscosity window & Approximately 10--500 mPa s tested; 500 mPa s retained sliding response better than low-viscosity comparison under reported scouring & Epoxy--TiO$_2$ SLIPS with abrasion/scouring & Study-specific design window; no universal optimum & Yang et al. \\
Lubricant-infused immersion & 40 days artificial seawater; roughness about $1.3\pm0.2~\mu$m for one structured morphology & Cerium-stearate structured SLIPS & Nanostructure improved retention relative to comparison morphology; loss rate NR & Yang et al. \\
Biolubricant infusion & Swelling ratio $1.3\pm0.1$ untreated methyl-oleate PDMS vs $1.9\pm0.4$ after UV pretreatment; mussel test two weeks & PDMS infused with methyl oleate & Infusion/short-term antifouling evidence; long-term lubricant-loss rate NR & Basu et al. \\
\bottomrule
\end{tabularx}
\end{table*}

\begin{figure}[t]
\centering
\fbox{\begin{minipage}{0.95\columnwidth}
\centering
\textbf{Biofouling formation} $\rightarrow$ conditioning film $\rightarrow$ biofilm $\rightarrow$ macrofouling\\[2pt]
$\downarrow$\\[-2pt]
\textbf{Surface response}: chemistry, roughness, wettability, hydration, lubrication\\[2pt]
$\downarrow$\\[-2pt]
\textbf{Control choice}: inhibit settlement, reduce adhesion, release, or remove\\[2pt]
$\downarrow$\\[-2pt]
\textbf{Service condition}: immersion, motion, abrasion, ageing, cleaning\\[2pt]
$\downarrow$\\[-2pt]
\textbf{Engineering outcomes}: roughness, drag, powering, fuel, maintenance\\[2pt]
$\downarrow$\\[-2pt]
\textbf{Sustainability outcomes}: release, ecotoxicity, cost, life-cycle burden
\end{minipage}}
\caption{Conceptual framework linking fouling formation, surface response, maintenance, hydrodynamics, and sustainability. The arrows are a synthesis of the reviewed evidence, not a single experimental pathway.}
\label{fig:framework}
\end{figure}

\begin{figure}[t]
\centering
\fbox{\begin{minipage}{0.95\columnwidth}
\centering
\textbf{Sustainable control portfolio}\\[4pt]
\begin{tabularx}{0.95\columnwidth}{@{}>{\centering\arraybackslash}X >{\centering\arraybackslash}X >{\centering\arraybackslash}X@{}}
Passive & Hybrid & Active\\
\shortstack{Silicone / hydrated\\low-adhesion surfaces} & \shortstack{Nanocomposite\\photocatalytic or\\bio-derived coating} & \shortstack{Grooming, waterjet,\\ultrasound, robot,\\laser, UV}\\[6pt]
\multicolumn{3}{>{\centering\arraybackslash}p{0.95\columnwidth}}{\textbf{Asset and operating profile: hull | net | membrane | submerged structure}}\\[6pt]
\multicolumn{3}{>{\centering\arraybackslash}p{0.95\columnwidth}}{\textbf{Decision criteria: fouling target | durability | drag | release | cost | waste}}\\
\end{tabularx}
\end{minipage}}
\caption{Conceptual taxonomy of control technologies and the common decision criteria needed for comparison.}
\label{fig:taxonomy}
\end{figure}

\begin{figure*}[!t]
\centering
\footnotesize
\renewcommand{\arraystretch}{1.1}
\fbox{\begin{minipage}{0.95\textwidth}
\centering
\textbf{Identification}\\
Scopus, 2015--2026, searched 5 Sept. 2026\\
8 tiered Boolean queries: 21{,}096 papers / 2{,}447 reviews (raw)
\end{minipage}}\\[2pt]
$\big\Downarrow$\\[2pt]
\fbox{\begin{minipage}{0.95\textwidth}
\centering
\textbf{Screening (title-level)}\\
1{,}160 unique records exported after per-tier fetch and de-duplication\\
Tier-1 strict niche: 200 screened
\end{minipage}}\\[2pt]
$\big\Downarrow$ \textit{\footnotesize 199 excluded (false positives, 99.5\%)}\\[2pt]
\fbox{\begin{minipage}{0.95\textwidth}
\centering
\textbf{Eligibility}\\
Tier-1 broadened: 166 screened\\
Adjacent Tier-2 pillars: 378 screened (contextual, not niche-specific)\\
Existing reviews on the strict niche: 69 screened, 0 genuine
\end{minipage}}\\[2pt]
$\big\Downarrow$\\[2pt]
\fbox{\begin{minipage}{0.95\textwidth}
\centering
\textbf{Included (strict niche)}\\
\textbf{1} title-verified paper\\
Verdict: intersection immature $\rightarrow$ retain broad scope
\end{minipage}}
\caption{PRISMA-style selection funnel for the Scopus feasibility search on the biomass/pyrogenic-carbon $\times$ marine-antifouling intersection (Tier 1). This funnel is separate from the 81-study curated corpus (Table~\ref{tab:taxonomy}) that underlies the main synthesis; detailed search counts are provided in the Supplementary Material.}
\label{fig:prisma_flow}
\end{figure*}

